% Generated draft: requires mathematical and source review.
\documentclass[11pt]{article}
\usepackage[margin=25mm]{geometry}
\usepackage{fontspec}
\setmainfont{Latin Modern Roman}
\usepackage{amsmath,amssymb,mathtools,unicode-math}
\setmathfont{Latin Modern Math}
\usepackage{enumitem,needspace,xurl,hyperref}
\setlength{\parindent}{0pt}
\setlength{\parskip}{5pt}
\setlength{\emergencystretch}{3em}
\newcommand{\N}{\mathbb{N}}
\newcommand{\Z}{\mathbb{Z}}
\newcommand{\Q}{\mathbb{Q}}
\newcommand{\R}{\mathbb{R}}
\newcommand{\C}{\mathbb{C}}
\newcommand{\F}{\mathbb{F}}
\newcommand{\A}{\mathbb{A}}
\newcommand{\PP}{\mathbb{P}}
\newcommand{\E}{\mathbb{E}}
\newcommand{\eps}{\varepsilon}
\newcommand{\dd}{\,\mathrm{d}}
\begin{document}
\section*{UnsolvedMath ID: 7200062}
\textbf{Source problem code:} catalogue entry
\subsection{Definitions and mathematical statement}
Let \(P\subset\R^3\) be a convex polyhedron. In the usual edge-net convention, choose uncut face-adjacency edges forming a tree, so that its rigid polygonal faces can be flattened to a connected planar development \(N\) with pairwise disjoint face interiors. Cut edges may separate, while uncut edges remain hinges. This describes a given valid net; the question is not merely whether some net exists.

A customary blooming is a continuous hinge motion from the original polyhedron to that planar net: each face remains congruent to itself, shared hinge edges stay identified, collisions are avoided, and each retained interior dihedral angle opens monotonically from its initial value to \(\pi\). In this convention the conjecture is
\[
\forall P\text{ convex polyhedron}\ \forall N\text{ an admissible net of }P,
\qquad \exists\text{ a blooming from }P\text{ to }N.
\]
The motion may involve many hinges simultaneously; no specified opening order or single-hinge-at-a-time requirement is imposed. It must reach the prescribed net, not a different convenient unfolding.

The extracted statement does not define net or blooming. Whether arbitrary cuts through faces are allowed, whether blooming requires monotone angles, and how boundary contacts are treated must be checked against the intended convention. The edge-net/monotone formulation above is made explicit for review rather than silently identifying it with every collision-free unfolding notion.

The supplied research summary reports examples/families admitting bloomings and no verified general resolution. A failed particular local opening strategy is not a counterexample unless every admissible blooming to that net is excluded. Its reported 2024--2026 search is metadata, not an independent status audit here.
\subsection{Short English statement}
Can every prescribed nonoverlapping net of a convex polyhedron be reached by a continuous collision-free blooming, in the intended monotone-hinge sense? The quantifier is over every valid net, not just one net per polyhedron.
\subsection{Sources}
\begin{itemize}
\item Local UnsolvedMath registry, permanent ID \(7200062\), exact source code \texttt{catalogue entry}, ``Connelly's blooming conjecture''; attribution to Connelly is supplied by the title and statement.
\item Wikipedia, \emph{List of unsolved problems in mathematics}, pinned revision 1366636767, Geometry bullet 62; \url{https://en.wikipedia.org/w/index.php?title=List_of_unsolved_problems_in_mathematics&oldid=1366636767#Geometry}. MediaWiki-source extraction; recorded access 2026-07-29.
\item Supplied OPEN-TRIAGE research summary mentions some Platonic and other polyhedron examples without precise references. No primary Connelly citation or complete general-resolution audit is supplied.
\end{itemize}
\end{document}
